\chapter{Concluding chapter}
\label{ch:conclusion}

We used observational data to understand how longer and warmer seqaosns change growth for juvenile and mature trees. Our findings are significant because they suggest that temperate forest recruitement may benefit from longer and warmer seasons which could even further be beneficial for range expansion. Interestingly and in contrast to our expectation, growth provenance did not influence growth, suggesting that the emphasis on local adaptation may not be as important as previously sugested. In addition mature trees store most of the forests living biomass, but they failed to grow more with with longer growing seasons. Only one was negatively affected by warmer seasons, but they others did not change their groth. This absence of trend could be the result of a poorly understood multi-year lag effect of carbon allocation to storage instead of growth. Taken together, these findings reveal how the complex nature of climate change, when coupled with the also complex dynamic of tree growth and looking at these across many species may be much harder to answer clearly how growth responds to longer and warmer seasons. While observational data is useful, it has limits. This is why experiments are fundamental to tease apart the growth drivers across multiple species.

During my MSc, I also investigated the impact of a growing season extension and different stressors on growth across 10 species (Figure \ref{fig:effects}). First, I participated in an experiment that looked at the effects of drought and defoliation (mimicing insect outbreaks) intensity and timing on the phenology and growth (Figure \ref{fig:effects}: Phaenoflex). This experiment will provide valuable insights into the often proposed mechanisms that are likely to increase in frequency with climate change and that can negatively impact tree growth potential to increase with longer and warmer seasons. In parralel to my two chapters, I also conducted a large-scale experiment in which we subjected seven species of juvenile trees to spring and fall treatments in a full factorial design. This experiment will be particularly helpful in confirming whether or not our observed reponses of the trees in a non-experimental design were correct and if they are not, we will be able to tease apart the effect of each season on growth with greater confience. I terminated this experiment last autumn, I still need to finish cleaning and analyse the data as well as writing this paper, which will be integrated as my first chapter in my PhD thesis that will start in January 2027.